\chapter{CPU decode 实战：从标量 GEMV 到 packed 低比特微内核}

\section{decode 优化先看形态}

7B 推理的 decode 阶段每次只处理一个新 token。对线性层来说，输入是一个 hidden 向量，权重是一个大矩阵，输出是一个新向量。这更像 GEMV，而不是大 batch GEMM。CPU 优化要先承认这个形态：每 token 反复流过大块权重，权重字节和 KV 字节很容易支配延迟。

\begin{keyidea}
CPU decode 微内核的目标不是写一个“通用最快矩阵乘”，而是为单 token 或小 batch 的固定权重读流建立可验证的 packed layout、累加策略、tail path、线程切分和性能账本。
\end{keyidea}

本章沿同一个 linear 语义推进：标量 fp32 reference，普通 int8 baseline，packed int8，int4 packed，SIMD 微内核，线程切分，NUMA 和 profiler。每一步都要能回到 oracle。

\section{标量 reference 不是性能 baseline}

reference 的职责是正确，不是代表可接受性能。

\begin{lstlisting}[numbers=none,caption={fp32 reference GEMV}]
for out in 0..O:
  acc = 0
  for k in 0..K:
    acc += weight[out, k] * input[k]
  output[out] = acc
\end{lstlisting}

性能 baseline 可以仍然是标量，但要固定测量条件：连续内存、关闭 debug 检查、固定线程数、固定输入、warmup、重复次数和输出校验。否则 benchmark 比较的是调试开销或缓存偶然性。

\section{int8 权重 baseline}

最简单的 int8 权重实现是边读边反量化：

\begin{lstlisting}[numbers=none,caption={int8 权重 + fp32 输入 baseline}]
for out in 0..O:
  acc = 0
  scale = scales[out]
  for k in 0..K:
    w = float(weight_i8[out, k]) * scale
    acc += w * input[k]
  output[out] = acc
\end{lstlisting}

这比 fp32 权重少读字节，但每个权重多了转换和 scale。若 \code{scale} 是 per-channel，它可以在外层读取一次；若是 per-group，内层每隔 group size 切换 scale。这里必须先有 QuantDesc：

\begin{lstlisting}[numbers=none,caption={int8 QuantDesc 最小字段}]
QuantDesc:
  bit_width = 8
  signed = true
  scale_policy = per_channel | per_group
  group_size
  zero_point_policy
  scale_dtype
  axis
\end{lstlisting}

没有 descriptor，kernel 无法知道 scale 怎样对应权重，也无法写通用 oracle。

\section{packed int8：先改变地址流}

原始权重按 \code{[out,k]} 行主序存储，对单输出行扫描友好。但 SIMD 微内核通常想一次计算多个 \code{out}，让同一个输入 \code{input[k]} 广播到多个输出 accumulator。为了让多个输出通道的权重连续，需要把小面板打包。

\begin{lstlisting}[numbers=none,caption={packed int8 panel 示意}]
for out_block in 0..O step BN:
  for k in 0..K:
    for oi in 0..BN:
      packed.append(weight_i8[out_block + oi, k])
\end{lstlisting}

打包后，内核的热循环读取 \code{packed[k, oi]} 是连续的：

\begin{lstlisting}[numbers=none,caption={packed GEMV 热循环}]
for out_block in 0..O step BN:
  acc[BN] = 0
  for k in 0..K:
    x = input[k]
    wv = load packed weights for BN outputs at this k
    acc += x * dequant(wv)
  store acc
\end{lstlisting}

packing 的成本应发生在模型 prepare 阶段，而不是每次请求 decode 阶段。plan report 要记录 packed weight 字节数、packing 时间和 layout id。否则上线后可能发现首个请求很慢，却不知道慢在 prepare。

\section{int4 packed：字节减少，解包增加}

int4 把两个 4-bit 权重放进一个 byte。常见 signed int4 范围是 \code{-8..7} 或通过 zero point 表示。解包需要取低四位和高四位，再做符号扩展或减 zero point。

\begin{lstlisting}[numbers=none,caption={int4 解包语义}]
byte b contains two values:
  lo = b & 0x0f
  hi = (b >> 4) & 0x0f

signed value:
  if nibble >= 8:
    value = nibble - 16
  else:
    value = nibble
\end{lstlisting}

int4 的优势是权重字节减半，风险是解包、scale metadata 和误差。group quant 下每 \code{G} 个权重共享一个 scale。若 \code{K} 不能整除 \code{G}，最后一组是 tail group。byte oracle 必须覆盖 nibble 顺序、符号扩展、group tail 和 scale 索引。

\section{累加器：低比特不等于低精度累加}

低比特权重不代表可以用低精度累加。常见策略有两类：

\begin{itemize}
  \item 权重低比特，输入 fp16/fp32，反量化后用 fp32 accumulator。
  \item 权重和激活都量化，用 int32 accumulator，再 requantize 或转回 fp32。
\end{itemize}

第一类实现简单，适合学习和 CPU baseline；第二类更接近完整整数推理，但需要处理激活 scale、zero point 修正、bias scale 和输出 requantize。第三册前面量化章节已经讲过 zero point 修正，CPU 微内核进入整数路径时必须复用同一合同。

\begin{warningbox}
不要把 int4 权重直接乘成 int8 accumulator。归约维度 \code{K=4096} 时，误差和溢出风险都不可接受。accumulator dtype 是 kernel descriptor 的核心字段。
\end{warningbox}

\section{SIMD 微内核的寄存器账本}

假设一次计算 \code{BN=16} 个输出通道。内核维护 16 个 accumulator lane。每个 \code{k}：

\begin{enumerate}
  \item 读取一个 \code{input[k]}。
  \item 广播到 SIMD 寄存器。
  \item 读取 packed 权重的 16 个值。
  \item 转换或解包到可乘格式。
  \item 乘加到 accumulator。
\end{enumerate}

\begin{lstlisting}[numbers=none,caption={抽象 SIMD GEMV 微内核}]
acc = zero_vector(BN)

for k in 0..K:
  x = broadcast(input[k])
  w = load_packed_weight_vector(k)
  wf = convert_or_dequant(w, scale_for(k))
  acc = fma(x, wf, acc)

store output block
\end{lstlisting}

这个账本要和 ISA 绑定。AVX2、AVX-512、NEON、SVE 的向量宽度、转换指令、整数 dot 能力和 mask tail 都不同。因此代码结构应分成 capability、descriptor、kernel 实现和 fallback，不要在业务调用点写一堆 ISA 分支。

\section{tail path：最后一块经常出错}

若输出通道 \code{O} 不是 \code{BN} 的倍数，最后一个 block 需要 tail。常见策略：

\begin{itemize}
  \item 单独标量处理 tail，简单但慢一点。
  \item SIMD mask store，要求 ISA 支持。
  \item padding packed weight 到 \code{BN} 对齐，输出只写有效部分。
\end{itemize}

若归约维度 \code{K} 不是 group size、解包块或向量宽度的倍数，也需要 tail。对 int4，\code{K} 为奇数还会出现半个 byte 的语义问题。虽然真实 7B 的很多维度规整，测试必须包含不规整维度，因为模型变体、lm head、adapter 或小 fixture 都可能触发。

\section{线程切分：按输出通道还是按层}

decode 的单个 linear 可以按输出通道切分：不同线程计算不同 \code{out} block，共享同一 input，读取不同权重。优点是没有输出写冲突；缺点是每个线程都读一遍 input，且小层线程开销可能过大。

\begin{lstlisting}[numbers=none,caption={按输出通道切分}]
parallel_for out_block_range:
  run_packed_gemv(input, packed_weight_range, output_range)
\end{lstlisting}

也可以在线路更高层调度：不同层或不同算子顺序依赖强，不适合随意并行；不同 request 的 decode step 可以并行，但会增加权重竞争和 cache 压力。最稳的起点是单算子内按输出 block 切分，并记录线程数和分区。

\begin{center}
\begin{tabular}{p{0.22\linewidth}p{0.34\linewidth}p{0.28\linewidth}}
\toprule
切分方式 & 优点 & 风险 \\
\midrule
按输出通道 & 简单，无输出冲突 & 权重带宽争用、线程开销 \\
按 batch/session & 可提高吞吐 & 状态复杂，p95 可能变差 \\
按层流水 & 理论并行 & decoder 层依赖强，收益有限 \\
\bottomrule
\end{tabular}
\end{center}

\section{NUMA 和权重放置}

多路 CPU 或 NUMA 机器上，权重放在哪个 NUMA node 会影响带宽。若线程在 node 0，权重页却在 node 1，远端内存访问会增加延迟。最小工程策略：

\begin{itemize}
  \item prepare 阶段按计划触碰 packed weight，让页分配到目标 node。
  \item 线程池固定亲和，避免跨 node 随机迁移。
  \item plan report 记录线程数、亲和策略和 NUMA 策略。
  \item benchmark 报告硬件拓扑，不把不同机器数字混比。
\end{itemize}

NUMA 优化不要太早做，但 trace 要保留足够字段。否则后期发现 p95 抖动，只能猜测是调度还是内存放置。

\section{perf counters：用证据判断瓶颈}

CPU 优化不能只看耗时。至少要观察：

\begin{itemize}
  \item cycles 和 instructions，判断 IPC。
  \item cache miss 和 LLC miss，判断权重/KV 读流。
  \item branch misses，判断 tail 或分支复杂度。
  \item vector instruction 比例，确认 SIMD path 真的执行。
  \item memory bandwidth，判断是否接近带宽上限。
\end{itemize}

如果引入 SIMD 后 instructions 降了但 LLC miss 上升，可能 packing layout 破坏了局部性。如果线程数增加后 tokens/s 不升反降，可能已经带宽饱和或 NUMA 远端访问增加。如果 int4 比 int8 不快，可能解包和 scale 成本抵消了字节收益。

\section{微内核进入主线的验收}

一个 CPU decode 微内核进入主线前，至少要满足：

\begin{rubricbox}
\begin{itemize}
  \item descriptor 明确 dtype、layout、group size、accumulator、tail policy 和 ISA。
  \item packed weight oracle 能把 packed bytes 映射回原始权重。
  \item operator oracle 覆盖整齐尺寸、输出 tail、归约 tail、非对齐输入和 group tail。
  \item layer oracle 覆盖至少一个 tiny decoder block。
  \item performance report 分开 pack time、kernel time、decode step time 和内存峰值。
  \item unsupported ISA 有明确 fallback 或 plan 阶段拒绝。
  \item 反汇编或 profiler 证明确实走到目标 SIMD path。
\end{itemize}
\end{rubricbox}

这条验收线比“benchmark 快了”严格，但它能保护项目长期演进。否则下一次改量化、改 layout、接 GPU fallback 时，很容易破坏 CPU 路径而不自知。

\section{本章验收线}

读完本章，应该能把 CPU decode 优化拆成可实现任务：

\begin{itemize}
  \item 从 fp32 reference 建立正确性基准，再写 int8/int4 baseline。
  \item 理解 packed layout 改变的是多个输出通道的权重地址流。
  \item 知道 int4 的 nibble 顺序、符号扩展、group scale 和 tail 为什么必须测试。
  \item 能说明 accumulator dtype 为什么是 kernel descriptor 的核心字段。
  \item 能把 SIMD 微内核写成 input broadcast、packed load、dequant、FMA、store 的账本。
  \item 能设计线程切分、NUMA 策略、perf counters 和主线验收。
\end{itemize}
